\scanpage{100}

Cela résulte en effet \inserted{dans le cas où $N'\not\simeq N_\infty$, $N''\not\simeq N_\infty$,} du deuxième th de convexité pour opérateurs, et des formules $\rho_1(uv)\ldots\rho_n(uv)\leq(\rho_1(u)\rho_1(v))\ldots(\rho_n(u)\rho_n(v))$ (N\textsuperscript{o} \unsure{,} formule (37)). -- En particulier, \deleted{si $N=$ \unsure{la}} pour \deleted{\unsure{N}} deux normes polaires $N$ et $N^\circ$, on trouve

% slanted
\margin{Cas $N'\simeq N_\infty$ ou $N''\simeq N_\infty$\,: si l'on \unsure{adm.} ceci $u,v$ \uline{cpt} on s'y ramène facilement dans le cas général, car en tous cas $N(u)=\lim_P N(\unsure{P_n u|Pu})$ $P$ projecteurs orthogonaux de rang fini.}

\[
\boxed{\lVert uv\rVert_1\leq N(u)N^\circ(v)}\tag*{\textcircled{17}}
\]

On trouve aussi \deleted{direct$^{\mathrm{mt}}$} \deleted{\unsure{On a obtenu aussi} $\boxed{N(uv)\leq\lVert u\rVert N(v),\quad N(uv)\leq\lVert v\rVert N(u)}$}

Le dual de $L_N(E,F)$ s'identifie à un sous-espace du dual de $L_{N_1}(E,F)=L^{\unsure{1}}(E,F)$, donc de $L(F,E)$. Ce sous-espace contient $L^{N^\circ}(F,E)$, et l'application $L^{N^\circ}(F,E)\to(L_N(E,F))'$ est un isomorphisme \uline{métrique}.

Montrons que c'est \uline{sur}. On peut supposer que $N\not\simeq N_1$, prouvons que \deleted{\unsure{si}} \inserted{si un} $\Phi\in(L_N(E,F))'$, définit \inserted{$v\in L(F,E)$, \uline{on} a $v\in L^{N^\circ}(F,E)$. Or la norme de} \deleted{un op. \uline{compact} de $F$ dans $E$. \unsure{Comme} $v^*v$ \unsure{la dual} utilise $|\operatorname{Tr}vu|\leq M N(u)$ pour $u\in E'\otimes F$, or $w=v^*v\in L(E,E)$ aura la propriété analogue, prouvons que $w$ est}

$v=U|v|$, il suffit de prouver que $|v|\in L^{N^\circ}(E)$\unsure{,|$F$|$E$} or $|v|=U'v$, donc a la \unsure{$\subset$|même} propriété analogue à $v$ ($|\operatorname{Tr}(vu)|\leq M N(u)$ pour $u\in E'\otimes F$, $M$ constante). On est ramené à $F=E$, $v\geq 0$.

$v$ est cpt\,: sinon il existerait un op. continu $w$ tel que \deleted{\unsure{$v$}} $wv=P$, projecteur hermitien de rang infini, \deleted{\unsure{on aurait la}} on aurait encore $w\in(L_N(E,E))'$, ce qui est impossible, \deleted{\unsure{on suppose un projecteur}} \inserted{(en prenant} une base orthonormale dans $P$, \deleted{\unsure{en aurait}}\unsure{on voit qu'on|on soit qu'on} aurait \deleted{$\operatorname{Tr}\unsure{p_n}P=$} $1\in(\ell_N)'=\ell^{N^\circ}$, ce qui n'est pas. Donc

on aura $v=\sum\rho_i\bar e_i\otimes e_i$, et sur \uline{la \unsure{base|même|dite|trouve} formule} on voit que $(\rho_i)\in(\ell_N)'=\ell^{N^\circ}$, d'où la conclusion. -- Comme conséquence, on voit donc que $L^{N^\circ}(E,F)$ est \uline{toujours} complet, donc tout espace $L^N(E,F)$ est complet donc aussi $L_N(E,F)$

\margin{?}
% archive sheet number: 101
